\icode{E-MK-MU-L2-260928-002}
（多选）
广州正规划与其他超大城市间高速磁悬浮通道及实验线建设，预留京港澳、沪广高速磁悬浮通道。某磁悬浮列车进站过程中的 \(\dfrac{x}{t} - t\) 图像如图所示，进站过程的运动可视为匀减速直线运动。下列说法正确的是\paren[BC]

\begin{tikzpicture}[>=stealth,font=\scriptsize]
	\draw[->,blue] (0,0)node[below left]{\(O\)}--(3,0)node[below]{\(t\)};
	\draw[->,blue] (0,0)--(0,3)node[right]{\(\frac{x}{t}\)};
	\draw[dashed,cyan] (0,1.3)--(1,1.3)--(1,0);
	\draw[purple] (0,2.6)node[left,blue]{\(b\)}--(2,0)node[below,blue]{\(t_0\)};
	\node[blue] at(0,1.3)[left]{\(\frac{b}{2}\)};
	\node[blue] at(1,0)[below]{\(\frac{t_0}{2}\)};
\end{tikzpicture}
\begin{choices}
	\item 列车进站的初速度大小为 \(2b\)
	\item 列车进站过程中的加速度大小为 \(\frac{2b}{t_0}\)
	\item 列车从开始进站到停止所用的时间为 \(\frac{t_0}{2}\)
	\item 列车进站过程通过的距离为 \(\frac{bt_0}{2}\)
\end{choices}
\begin{answer}
	BC\qquad 提示：对匀变速直线运动\(\frac{x}{t} = v_0 + \frac{1}{2}at\)
\end{answer}